% problem_id: S001-lemma-quasi-finite-over-base
% source_id: S001 (Stacks Project)
% source_file: more-groupoids.tex
% statement_locator: more-groupoids.tex lines 2178--2228
% proof_locator: more-groupoids.tex lines 2230--2344
% env: lemma
% id_note: label 在本来源内唯一
% pairing: tight (verified)
% status: complete (statement + author proof verbatim + reason + locators)
%
\begin{lemma}
\label{lemma-quasi-finite-over-base}
Let $S$ be a scheme.
Let $(U, R, s, t, c)$ be a groupoid scheme over $S$.
Let $p \in S$ be a point, and let $u \in U$ be a point lying over $p$.
Assume that
\begin{enumerate}
\item $U \to S$ is locally of finite type,
\item $U \to S$ is quasi-finite at $u$,
\item $U \to S$ is separated,
\item $R \to S$ is separated,
\item $s$, $t$ are flat and locally of finite presentation, and
\item $s^{-1}(\{u\})$ is finite.
\end{enumerate}
Then there exists an \'etale neighbourhood $(S', p') \to (S, p)$ with
$\kappa(p) = \kappa(p')$ and a base change diagram
$$
\xymatrix{
R' \amalg W'
\ar@{=}[r] &
S' \times_S R
\ar[r] \ar@<2ex>[d]^{s'} \ar@<-2ex>[d]_{t'} &
R \ar@<1ex>[d]^s \ar@<-1ex>[d]_t \\
U' \amalg W
\ar@{=}[r] &
S' \times_S U
\ar[r] \ar[d] &
U \ar[d] \\
 &
S' \ar[r] &
S
}
$$
where the equal signs are decompositions into open and closed
subschemes such that
\begin{enumerate}
\item[(a)] there exists a point $u'$ of $U'$ mapping to $u$ in $U$,
\item[(b)] the fibre $(U')_{p'}$ equals $t'\big((s')^{-1}(\{u'\})\big)$
set theoretically,
\item[(c)] the fibre $(R')_{p'}$ equals $(s')^{-1}\big((U')_{p'}\big)$
set theoretically,
\item[(d)] the schemes $U'$ and $R'$ are finite over $S'$,
\item[(e)] we have $s'(R') \subset U'$ and $t'(R') \subset U'$,
\item[(f)] we have
$c'(R' \times_{s', U', t'} R') \subset R'$
where $c'$ is the base change of $c$, and
\item[(g)] the morphisms $s', t', c'$ determine a groupoid structure
by taking the system
$(U', R', s'|_{R'}, t'|_{R'}, c'|_{R' \times_{s', U', t'} R'})$.
\end{enumerate}
\end{lemma}

\begin{proof}
Let us denote $f : U \to S$ the structure morphism of $U$.
By assumption (6) we can write $s^{-1}(\{u\}) = \{r_1, \ldots, r_n\}$.
Since this set is finite, we see that $s$ is quasi-finite at each of
these finitely many inverse images, see
Morphisms, Lemma \ref{morphisms-lemma-finite-fibre}.
Hence we see that $f \circ s : R \to S$ is quasi-finite at each $r_i$
(Morphisms, Lemma \ref{morphisms-lemma-composition-quasi-finite}).
Hence $r_i$ is isolated in the fibre $R_p$, see
Morphisms, Lemma \ref{morphisms-lemma-quasi-finite-at-point-characterize}.
Write $t(\{r_1, \ldots, r_n\}) = \{u_1, \ldots, u_m\}$.
Note that it may happen that $m < n$ and note that
$u \in \{u_1, \ldots, u_m\}$.
Since $t$ is flat and locally of finite presentation,
the morphism of fibres $t_p : R_p \to U_p$ is flat and locally of
finite presentation (Morphisms,
Lemmas \ref{morphisms-lemma-base-change-flat} and
\ref{morphisms-lemma-base-change-finite-presentation}),
hence open (Morphisms,
Lemma \ref{morphisms-lemma-fppf-open}).
The fact that each $r_i$ is isolated in $R_p$ implies that
each $u_j = t(r_i)$ is isolated in $U_p$. Using
Morphisms, Lemma \ref{morphisms-lemma-quasi-finite-at-point-characterize}
again, we see that $f$ is quasi-finite at $u_1, \ldots, u_m$.

\medskip\noindent
Denote $F_u = s^{-1}(u)$ and $F_{u_j} = s^{-1}(u_j)$ the scheme theoretic
fibres. Note that $F_u$ is finite over $\kappa(u)$ as it is locally of finite
type over $\kappa(u)$ with finitely many points (for example it follows from
the much more general
Morphisms, Lemma
\ref{morphisms-lemma-locally-quasi-finite-qc-source-universally-bounded}).
By
Lemma \ref{lemma-two-fibres}
we see that $F_u$ and $F_{u_j}$ become isomorphic over a common
field extension of $\kappa(u)$ and $\kappa(u_j)$. Hence we see
that $F_{u_j}$ is finite over $\kappa(u_j)$. In particular we see
$s^{-1}(\{u_j\})$ is a finite set for each $j = 1, \ldots, m$.
Thus we see that assumptions (2) and (6) hold for each $u_j$ also
(above we saw that $U \to S$ is quasi-finite at $u_j$).
Hence the argument of the first paragraph applies to each $u_j$
and we see that $R \to U$ is quasi-finite at each of the points of
$$
\{r_1, \ldots, r_N\} = s^{-1}(\{u_1, \ldots, u_m\})
$$
Note that $t(\{r_1, \ldots, r_N\}) = \{u_1, \ldots, u_m\}$ and
$t^{-1}(\{u_1, \ldots, u_m\}) = \{r_1, \ldots, r_N\}$
since $R$ is a groupoid\footnote{Explanation in groupoid language:
The original set $\{r_1, \ldots, r_n\}$ was the set of arrows with
source $u$. The set $\{u_1, \ldots, u_m\}$ was the set of objects
isomorphic to $u$. And $\{r_1, \ldots, r_N\}$ is the set of all arrows
between all the objects equivalent to $u$.}. Moreover, we have
$\text{pr}_0(c^{-1}(\{r_1, \ldots, r_N\})) = \{r_1, \ldots, r_N\}$
and
$\text{pr}_1(c^{-1}(\{r_1, \ldots, r_N\})) = \{r_1, \ldots, r_N\}$.
Similarly we get $e(\{u_1, \ldots, u_m\}) \subset \{r_1, \ldots, r_N\}$
and $i(\{r_1, \ldots, r_N\}) = \{r_1, \ldots, r_N\}$.

\medskip\noindent
We may apply
More on Morphisms,
Lemma \ref{more-morphisms-lemma-etale-splits-off-quasi-finite-part-technical}
to the pairs
$(U \to S, \{u_1, \ldots, u_m\})$ and
$(R \to S, \{r_1, \ldots, r_N\})$
to get an \'etale neighbourhood $(S', p') \to (S, p)$
which induces an identification $\kappa(p) = \kappa(p')$
such that $S' \times_S U$ and $S' \times_S R$ decompose as
$$
S' \times_S U = U' \amalg W, \quad
S' \times_S R = R' \amalg W'
$$
with $U' \to S'$ finite and $(U')_{p'}$ mapping bijectively to
$\{u_1, \ldots, u_m\}$, and $R' \to S'$ finite and
$(R')_{p'}$ mapping bijectively to $\{r_1, \ldots, r_N\}$.
Moreover, no point of $W_{p'}$ (resp.\ $(W')_{p'}$) maps to
any of the points $u_j$ (resp.\ $r_i$). At this point (a), (b), (c), and (d)
of the lemma are satisfied. Moreover, the inclusions of (e) and (f) hold
on fibres over $p'$, i.e., $s'((R')_{p'}) \subset (U')_{p'}$,
$t'((R')_{p'}) \subset (U')_{p'}$, and
$c'((R' \times_{s', U', t'} R')_{p'}) \subset (R')_{p'}$.

\medskip\noindent
We claim that we can replace $S'$ by a Zariski open neighbourhood
of $p'$ so that the inclusions of (e) and (f) hold.
For example, consider the set $E = (s'|_{R'})^{-1}(W)$.
This is open and closed in $R'$ and does not contain any points
of $R'$ lying over $p'$. Since $R' \to S'$ is closed,
after replacing $S'$ by $S' \setminus (R' \to S')(E)$ we reach a
situation where $E$ is empty. In other words $s'$ maps $R'$ into $U'$.
Note that this property is preserved under further shrinking $S'$.
Similarly, we can arrange it so that $t'$ maps $R'$ into $U'$.
At this point (e) holds. In the same manner, consider the set
$E = (c'|_{R' \times_{s', U', t'} R'})^{-1}(W')$.
It is open and closed in the scheme $R' \times_{s', U', t'} R'$
which is finite over $S'$, and does not contain any points lying
over $p'$. Hence after replacing $S'$ by
$S' \setminus (R' \times_{s', U', t'} R' \to S')(E)$
we reach a situation where $E$ is empty. In other words we obtain
the inclusion in (f). We may repeat the argument also with the identity
$e' : S' \times_S U \to S' \times_S R$ and the inverse
$i' : S' \times_S R \to S' \times_S R$ so that we may assume
(after shrinking $S'$ some more) that $(e'|_{U'})^{-1}(W') = \emptyset$
and $(i'|_{R'})^{-1}(W') = \emptyset$.

\medskip\noindent
At this point we see that we may consider the structure
$$
(U', R', s'|_{R'}, t'|_{R'}, c'|_{R' \times_{t', U', s'} R'},
e'|_{U'}, i'|_{R'}).
$$
The axioms of a groupoid scheme over $S'$ hold
because they hold for the groupoid scheme
$(S' \times_S U, S' \times_S R, s', t', c', e', i')$.
\end{proof}
